\begin{thebibliography}{99}
\raggedright

\bibitem{Fantechi}
Barbara Fantechi, \emph{Stacks for Everybody}, in
\emph{European Congress of Mathematics}, Barcelona 2000, Vol.~I,
Progress in Mathematics 201, Birkh\"auser, 2001, pp.~349--359.
\href{https://doi.org/10.1007/978-3-0348-8268-2_20}{DOI and published chapter}.
Local transcription: \path{../../Papers/Stacks/StacksForEverybody.tex}.
The source-PDF page markers in that file are used for its internal
pagination. Sections 3.2, 3.5, 4.1--4.3, 5, and 6.3 are used here.

\bibitem{Vistoli}
Angelo Vistoli, \emph{Notes on Grothendieck topologies, fibered
categories and descent theory}, author version of 2 October 2008.
\href{https://homepage.sns.it/vistoli/descent.pdf}{Author PDF};
\href{https://arxiv.org/abs/math/0412512}{arXiv record}.
The notes form the descent contribution to \emph{Fundamental
Algebraic Geometry: Grothendieck's FGA Explained}, AMS, 2005.
Numbering and pages in these notes refer to the 2008 author PDF:
Propositions 1.15, 4.31, 4.37; Theorems 4.23 and 4.38.

\bibitem{SGA1}
Alexander Grothendieck et al., \emph{Rev\^etements \'etales et groupe
fondamental (SGA 1)}, Documents Math\'ematiques 3, Soci\'et\'e
Math\'ematique de France, 2003, recomposed and annotated edition
of Lecture Notes in Mathematics 224 (1971).
\href{https://arxiv.org/abs/math/0206203}{arXiv edition},
\href{https://arxiv.org/pdf/math/0206203}{PDF}.
Expos\'e VIII, Proposition 7.8 is on printed p.~174 of the
recomposed edition (PDF p.~190; original margin pp.~224--225).

\bibitem{Alper}
Jarod Alper, \emph{Stacks and Moduli}, version of 5 January 2026.
\href{https://sites.math.washington.edu/~jarod/moduli-versions/moduli-1-5-26.pdf}{Dated author PDF}.
Local copy: \path{../../Papers/Stacks/JarodAlperModuliNotes.pdf}.
The dated version is important: its stable-curve material is
Chapter 6, and theorem numbering differs from earlier versions.
Principal passages: Sections 6.3--6.5, pp.~225--252;
Theorems 2.1.2, 4.2.1, 4.7.1, and 4.8.7; Appendix C.

\bibitem{DM}
Pierre Deligne and David Mumford, \emph{The irreducibility of the
space of curves of given genus}, Publications Math\'ematiques
de l'IH\'ES \textbf{36} (1969), 75--109.
\href{https://www.numdam.org/item/PMIHES_1969__36__75_0/}{Numdam record};
\href{https://www.numdam.org/item/PMIHES_1969__36__75_0.pdf}{original PDF}.
Sections 1 and 5; especially Theorem (1.2), Lemma (1.4),
Proposition (5.1), and Theorem (5.2). Pages in the text are
journal pages.

\bibitem{Knudsen}
Finn F. Knudsen, \emph{The projectivity of the moduli space of
stable curves, II: The stacks $M_{g,n}$}, Mathematica Scandinavica
\textbf{52} (1983), 161--199.
\href{https://doi.org/10.7146/math.scand.a-12001}{DOI and journal article}.
Theorem 2.7, journal p.~179, is the pointed proper smooth stack
theorem. Its historical ``algebraic stack'' terminology is the
Deligne--Mumford terminology used here.

\bibitem{Edidin}
Dan Edidin, \emph{Notes on the construction of the moduli space
of curves}, in \emph{Recent Progress in Intersection Theory},
Trends in Mathematics, Birkh\"auser, 2000, pp.~85--113.
\href{https://arxiv.org/abs/math/9805101}{arXiv record and text}.
Section 3 of the preprint, Theorems 3.2--3.4, pp.~15--17,
is a later exposition of the unpointed construction and reduction
arguments; it is supplementary rather than the main proof here.

\bibitem{FP}
William Fulton and Rahul Pandharipande, \emph{Notes on stable
maps and quantum cohomology}, in \emph{Algebraic Geometry---Santa
Cruz 1995}, Proceedings of Symposia in Pure Mathematics 62,
Part 2, AMS, 1997, pp.~45--96.
\href{https://arxiv.org/abs/alg-geom/9608011}{arXiv record};
\href{https://arxiv.org/pdf/alg-geom/9608011}{PDF}.
Citations use the preprint's printed pagination: Sections 1--2,
pp.~10--14; Section 4.2, Propositions 5--6, pp.~19--21;
Section 5.1, Lemma 8, pp.~26--27; Section 5.2, pp.~27--30.

\bibitem{BM}
Kai Behrend and Yuri Manin, \emph{Stacks of stable maps and
Gromov--Witten invariants}, Duke Mathematical Journal
\textbf{85} (1996), no.~1, 1--60.
\href{https://arxiv.org/abs/alg-geom/9506023}{arXiv record};
\href{https://arxiv.org/pdf/alg-geom/9506023}{PDF}.
Local text: \path{../../Papers/BehrendManin/BehrendManin.tex}.
Citations use arXiv pagination: Definitions 2.1, 3.2, 3.4,
3.13; Theorem 3.14; Proposition 4.1, Lemma 4.2, and
Corollary 4.8, pp.~12--28.

\bibitem{BCM}
Luca Battistella, Francesca Carocci, and Cristina Manolache,
\emph{Virtual Classes for the Working Mathematician}, SIGMA
\textbf{16} (2020), 026, 38 pages.
\href{https://doi.org/10.3842/SIGMA.2020.026}{DOI};
\href{https://www.emis.de/journals/SIGMA/2020/026/sigma20-026.pdf}{journal PDF}.
Local text: \path{../../Papers/VCWM/sigma20-026.tex}.
Only the concrete singular stable-map example is used:
Section 8.2, Example 8.4, p.~27; Section 8.3, equation (8.2),
p.~29. The paper's later-purpose machinery is not used in our
stack proofs.

\bibitem{Stacks}
The Stacks Project Authors, \emph{The Stacks Project},
\href{https://stacks.math.columbia.edu}{project website}.
Tags are permanent identifiers and are linked at their uses.
Main technical checks: \MCtag{023Q} (morphism descent),
\MCtag{0E6R} (relative dualizing sheaf), \MCtag{0B5Y}
(curve ampleness), \MCtag{0D2S} (relative ampleness),
\MCtag{02L0}, \MCtag{02L1}, \MCtag{02L2} (descent of
morphism properties). The project fills precise technical
gaps here; Fantechi, Vistoli, Alper, and Fulton--Pandharipande
provide the main pedagogical route.

\bibitem{Vakil}
Ravi Vakil, \emph{The moduli space of curves and Gromov--Witten
theory}, in \emph{Enumerative Invariants in Algebraic Geometry
and String Theory}, Lecture Notes in Mathematics 1947,
Springer, 2008, pp.~143--198.
\href{https://arxiv.org/abs/math/0602347}{arXiv record};
\href{https://math.stanford.edu/~vakil/files/cimefeb26.pdf}{author preprint}.
Sections 2.4--2.7, preprint pp.~9--11, give the supplementary
compactification discussion considered here.

\bibitem{Zvonkine}
Dimitri Zvonkine, \emph{An introduction to moduli spaces of
curves and their intersection theory}, version of 2 November 2014.
\href{https://math.hse.ru/data/2014/11/13/1101420186/2013.04.03.pdf}{dated lecture-note PDF}.
Section 1.4, pp.~14--18, gives geometric intuition; its
introductory account deliberately does not develop the stack
foundations.

\bibitem{Belfiori}
Filippo Belfiori, \emph{On the Construction of Some Stacks of
Curves and Surfaces}, master's thesis, Universit\`a di Bologna,
academic year 2023--24.
\href{https://amslaurea.unibo.it/id/eprint/35200/1/Tesi_BelfioriFilippo.pdf}{institutional thesis PDF}.
Section 1.1, Definition 1.6, Remark 1.7, Theorem 1.12,
Remark 1.13, and Theorem 1.18, pp.~2--5.

\bibitem{Monavari}
Sergej Monavari, \emph{Local Picture of Twisted Curves with
an Introduction to the Theory of Deligne--Mumford Stacks},
ALGANT master's thesis, 2018.
\href{https://www.math.u-bordeaux.fr/~ybilu/algant/documents/theses/Monavari.pdf}{institutional thesis PDF}.
Sections 1.2--1.3, pp.~11--16, were considered for introductory
fibered-category and stack definitions.

\end{thebibliography}
